\section{Einleitung}
\label{sec:einleitung}

\subsection{Motivation}

Webanwendungen verarbeiten fortlaufend Daten aus Formularen,
Datenbanken, APIs und dem Browserkontext. Zwei zentrale
Sicherheitsgrenzen sind dabei die Trennung von Daten und ausführbarem
Code sowie die Trennung verschiedener Web-Origins. XSS verletzt die
erste Grenze: Ein eigentlich als Text gedachter Wert wird als HTML
oder JavaScript interpretiert. CORS steuert die zweite Grenze: Ein
Server kann dem Browser erlauben, Antworten auch fremden Origins
zugänglich zu machen.

Der Scheduler eignet sich als praxisnahes Untersuchungsobjekt, weil
er zahlreiche nutzerkontrollierbare Stammdaten und Kalenderansichten,
ein Vue-Frontend, eine Spring-Boot-API und Keycloak zur
Authentifizierung kombiniert. Ein Fehler an der Ausgabegrenze kann
dadurch im Browser eines eingeloggten Nutzers ausgeführt werden. Eine
zu breite Origin-Freigabe kann wiederum die durch die Same-Origin
Policy vorgesehene Lesesperre aufheben.

\subsection{Zielsetzung und Leitfragen}

Ziel des Projekts war nicht eine generische Darstellung der beiden
Themen, sondern eine konkrete Sicherheitsprüfung des
Scheduler-Projekts mit nachvollziehbaren Live-Demos und umsetzbaren
Verbesserungen. Der Bericht beantwortet vier Leitfragen:

\begin{enumerate}
  \item Welche XSS- und CORS-relevanten Datenflüsse und
    Konfigurationen existieren im Scheduler?
  \item Welche Schwachstellen oder riskanten Konfigurationen lassen
    sich tatsächlich belegen?
  \item Wie können Ursache, Wirkung und Behebung lokal demonstriert werden?
  \item Welche Maßnahmen sollten auf \texttt{capstone-main} dauerhaft
    umgesetzt und getestet werden?
\end{enumerate}

\subsection{Abgrenzung und Berichtsstand}

Die Prüfung bezog sich auf getrackte Quell- und Konfigurationsdateien
aus \texttt{api/}, \texttt{web/}, Docker, Nginx und dem
Keycloak-Realm-Export \texttt{scheduler.json}. Die konkrete
Rendering-Stelle in \texttt{vue-cal} wurde zusätzlich in der
installierten Bibliothek untersucht. Produktive Systeme wurden nicht
aktiv angegriffen.

Der Bericht beschreibt den Stand nach dem vollständigen Rollout des
XSS-Fixes auf \path{capstone-main}. Die Demo-Setups wurden dagegen in
einem separaten, von \path{capstone-main} abgeleiteten
Security-Demobranch umgesetzt. Absichtlich schwache
CORS-Einstellungen, Umschaltlogik und Demo-Testdaten gehören
ausschließlich zu diesem Demonstrationsstand und nicht zum produktiven Fix.
